\documentclass[11pt]{article}
\usepackage[margin=1in]{geometry}
\usepackage{setspace}
\usepackage{natbib}
\usepackage[T1]{fontenc}
\usepackage{lmodern}

\title{Causal Discovery in Time Series: Introduction and Literature Review}
\author{}
\date{}

\begin{document}
\maketitle
\onehalfspacing

\section{Introduction}
Causal discovery in time series seeks to infer directional data-generating relationships from temporal observations. This goal is central in econometrics, climate science, neuroscience, and systems biology, where experimental interventions are often costly or infeasible. In contrast to static causal discovery, temporal data include lagged dependencies that can be exploited for directionality, but these same dependencies also amplify challenges related to autocorrelation, nonstationarity, hidden confounding, and feedback.

The modern time-series causal discovery landscape is built on several intellectual threads. The first is predictive causality in econometrics, especially Granger's operational notion that a variable $X$ ``causes'' $Y$ if past values of $X$ improve prediction of $Y$ beyond past $Y$ and other conditioning variables \citep{granger1969}. This paradigm was developed in multivariate settings through vector autoregressive (VAR) modeling \citep{sims1980} and extended to nonstationary systems via cointegration and error-correction formulations \citep{engle1987}. The second thread is graphical causality, formalized through structural and probabilistic causal models \citep{pearl2000,spirtes2000}, which clarifies assumptions such as causal sufficiency and faithfulness. A third thread is information-theoretic directionality, where transfer entropy generalizes predictive influence beyond linear-Gaussian assumptions \citep{schreiber2000}. More recently, machine learning has enabled scalable structure learning through continuous optimization formulations (e.g., NOTEARS) \citep{zheng2018} and large-scale nonlinear time-series methods such as PCMCI \citep{runge2019}.

Despite major progress, robust causal discovery from observational time series remains unresolved in many realistic regimes. Three factors explain this gap. First, identifiability depends on assumptions that are domain-dependent and frequently violated in practice (e.g., stationarity, no latent confounders, correct lag order). Second, methods are often evaluated under heterogeneous simulation protocols and preprocessing decisions, weakening cross-study comparability. Third, the objective itself is multifaceted: one may seek predictive directional links, mechanistic intervention-relevant causes, or graph-level summaries, and these goals are not always aligned.

This review synthesizes the literature around those tensions. We first summarize foundational definitions and assumptions, then organize methods by inferential principle (predictive, graphical, information-theoretic, and optimization-based), and finally discuss benchmarks, evaluation metrics, and open methodological gaps.

\section{Literature Review}
\subsection{Foundational Definitions and Assumptions}
The Granger framework \citep{granger1969} remains the most widely deployed baseline for time-series causality. In its classical linear form, one compares predictive models with and without lagged predictors from a candidate cause and tests whether forecast error decreases significantly. In multivariate settings, VAR formulations \citep{sims1980} provide a practical parameterization in which conditional Granger relations can be tested via coefficient restrictions.

However, this framework is inherently assumption-sensitive. Reliable inference typically requires (i) an adequate lag embedding, (ii) weak stationarity or transformations that induce it, (iii) absence of severe latent confounding, and (iv) model class adequacy (often linear-Gaussian in classical implementations). For nonstationary but cointegrated systems, error-correction models become essential to avoid spurious causality claims \citep{engle1987}. Frequency-domain perspectives further characterize directional dependence and feedback \citep{geweke1982}, emphasizing that causal interpretation can vary across timescales.

Graphical causal modeling provides a complementary language for assumptions and identifiability. The works of \citet{spirtes2000} and \citet{pearl2000} formalize how conditional independences constrain causal structure and when equivalence classes rather than unique graphs are recoverable. For time-series applications, these ideas motivate conditioning strategies and clarify how hidden common causes can induce misleading apparent directions.

\subsection{Information-Theoretic and Non-Gaussian Approaches}
Information-theoretic approaches aim to detect directed dependence beyond linear models. Transfer entropy \citep{schreiber2000} quantifies directed information flow by comparing uncertainty in future $Y_t$ with and without past $X_{t-\tau}$, conditioned on relevant histories. This makes transfer entropy attractive in nonlinear settings, though estimator choice and finite-sample bias can strongly affect results.

Non-Gaussian identification results also expanded the landscape. In particular, structural VAR estimation leveraging non-Gaussianity \citep{hyvarinen2010} showed that distributional asymmetries can improve identifiability where covariance-based methods are insufficient. These approaches are especially relevant in domains where heavy-tailed or skewed innovations are expected.

\subsection{Constraint-Based and Large-Scale Time-Series Discovery}
Constraint-based causal discovery, rooted in conditional independence testing, was systematized for static settings in \citet{spirtes2000} and adapted for temporal regimes in subsequent methods. A major advance is PCMCI and related developments for high-dimensional nonlinear time series \citep{runge2019}, which combine condition-selection and testing steps to improve power and reduce false positives under limited samples.

These methods offer a practical compromise between exhaustive conditioning and statistical efficiency. They are particularly valuable when dense conditioning sets would otherwise collapse test power, but performance remains sensitive to test calibration, lag search space, and sample size relative to graph complexity.

\subsection{Optimization-Based Structure Learning and Temporal Extensions}
Optimization-based graph learning became prominent with NOTEARS \citep{zheng2018}, which encodes acyclicity as a smooth constraint and enables continuous optimization over weighted adjacency matrices. Although originally formulated for i.i.d.\ data and static DAGs, NOTEARS reshaped causal discovery by integrating graph constraints into differentiable objectives and inspired temporal and nonlinear variants.

The main appeal of this family is computational scalability and compatibility with modern optimization tooling. The main caveat is that objective design and regularization encode strong inductive biases; consequently, estimated graphs can be stable numerically yet fragile causally when assumptions mismatch the data-generating process.

\subsection{Evaluation Protocols, Benchmarks, and Persistent Gaps}
Across paradigms, evaluation commonly emphasizes edge-level classification metrics (precision, recall, F1, false positive rate) and graph-level distances such as structural Hamming distance. Yet benchmarking remains difficult because reported performance depends strongly on simulation design (noise model, nonlinearities, lag depth, graph density), preprocessing, and whether latent confounding or measurement noise are present.

Several persistent gaps recur across the literature:
\begin{enumerate}
  \item \textbf{Latent confounding and selection effects}: many methods still rely on causal sufficiency assumptions from static settings \citep{pearl2000,spirtes2000}.
  \item \textbf{Nonstationarity and regime shifts}: methods robust to time-varying mechanisms remain comparatively underdeveloped in standard toolchains.
  \item \textbf{Lag-order and window selection}: practical heuristics can dominate downstream graph quality.
  \item \textbf{Cycle handling and contemporaneous effects}: many frameworks impose acyclicity or separate instantaneous effects by modeling choice rather than identification.
  \item \textbf{Interpretability vs.\ predictive performance}: high-performing predictive directionality does not always map to intervention-relevant causation.
\end{enumerate}

Overall, the field has matured from linear predictive tests to richly structured graphical and optimization-based approaches, but robust causal claims in realistic multivariate time series still require explicit assumption auditing, sensitivity analyses, and benchmark designs that reflect domain-specific failure modes.

\bibliographystyle{plainnat}
\bibliography{references}

\end{document}
